% Korrektheit fuer die konstantenfreie Sprache von Band 48.
\section{Korrektheit}

Die Objektsprache bleibt \(\mathcal L_\in\).
In der Metatheorie duerfen die Regeln aus Band~1 benutzt werden,
waehrend hier ihre Wirkung auf \emph{codierte} Herleitungen untersucht
wird. Bei der All-Einfuehrung ist die Variable in keiner offenen
Annahme frei; beim klassischen Widerspruchsschluss wird die
angenommene Negation entladen. Die in Band~1 erklaerten
Abhaengigkeitsmengen werden unten ausdruecklich mitgefuehrt.

\FormulaDefDeltaK[Abkuerzungen in der codierten Objektsprache]
{\begin{aligned}
 p\lor q&\coloneq\neg(\neg p\land\neg q),\\
 p\to q&\coloneq\neg(p\land\neg q),\\
 p\leftrightarrow q&\coloneq(p\to q)\land(q\to p),\\
 \forall x\,p&\coloneq\neg\exists x\,\neg p.
\end{aligned}}
{B48LogicalAbbreviations}{
 \DeltaRow{Formeln und Variable}{p,q\in\mathcal F,\quad x\in\mathbb N}
}
Dies praezisiert die zuvor verwendeten ueblichen Abkuerzungen.
Ihre semantischen Klauseln muessen aus der bereits konstruierten
Erfuellungsrelation folgen.

\FormulaThmDeltaK[Semantik der abgeleiteten logischen Zeichen]
{\begin{aligned}
 (M,s\models p\lor q)&\leftrightarrow
       ((M,s\models p)\lor(M,s\models q)),\\
 (M,s\models p\to q)&\leftrightarrow
       ((M,s\models p)\to(M,s\models q)),\\
 (M,s\models p\leftrightarrow q)&\leftrightarrow
       ((M,s\models p)\leftrightarrow(M,s\models q)),\\
 (M,s\models\forall x\,p)&\leftrightarrow
       \forall a\in D\;(M,s[x\mapsto a]\models p).
\end{aligned}}
{B48DerivedSemanticClauses}{
 \DeltaRow{Struktur und Belegung}{\mathsf{Str}(M),\quad s\in A_D}
 \DeltaRow{Formeln und Variable}{p,q\in\mathcal F,\quad x\in\mathbb N}
}
In den folgenden Tabellen kuerzen \(P\) und \(Q\) ausschliesslich
die metasprachlichen Aussagen \(M,s\models p\) und \(M,s\models q\)
ab.
\begin{tabproofsplitwide}
 \proofpartwideindR[Disjunktion]{(M,s\models p\lor q)\leftrightarrow(P\lor Q)}
  {(M,s\models p\lor q)\leftrightarrow(P\lor Q)}[B48SatOrClause]
 \proofstepwidestar[]{(M,s\models p\lor q)\leftrightarrow
                          \neg(\neg P\land\neg Q)}
  {\FormulaRefAuto{B48LogicalAbbreviations}[def],
   \FormulaRefAuto{B48Satisfaction}[def]}
 \proofstepwidestar[]{\neg(\neg P\land\neg Q)
                          \leftrightarrow(\neg\neg P\lor\neg\neg Q)}
  {\FormulaRefAuto{\neg(P\land Q)\eqvdash\neg P\lor\neg Q}[thm]}
 \proofstepwidestar[]{(\neg\neg P\lor\neg\neg Q)\leftrightarrow(P\lor Q)}
  {\FormulaRefAuto{P\lor Q\eqvdash\neg\neg P\lor\neg\neg Q}[thm],
   \FormulaRefAuto{B48IffSymmetry}[thm]}
 \proofstepwidestar[]{(M,s\models p\lor q)\leftrightarrow(P\lor Q)}
  {\rLRS{3,\rLRS{2,1}}}
 \closeproofpartwideindR

 \proofpartwideindR[Implikation]{(M,s\models p\to q)\leftrightarrow(P\to Q)}
  {(M,s\models p\to q)\leftrightarrow(P\to Q)}[B48SatImpClause]
 \proofstepwidestar[]{(M,s\models p\to q)\leftrightarrow\neg(P\land\neg Q)}
  {\FormulaRefAuto{B48LogicalAbbreviations}[def],
   \FormulaRefAuto{B48Satisfaction}[def]}
 \proofstepwidestar[]{\neg(P\land\neg Q)\leftrightarrow(\neg P\lor\neg\neg Q)}
  {\FormulaRefAuto{\neg(P\land Q)\eqvdash\neg P\lor\neg Q}[thm]}
 \proofstepwidestar[]{Q\leftrightarrow\neg\neg Q}
  {\FormulaRefAuto{P\eqvdash\neg\neg P}[thm]}
 \proofstepwidestar[]{\neg(P\land\neg Q)\leftrightarrow(\neg P\lor Q)}
  {\rLRS{3,2}}
 \proofstepwidestar[]{(P\to Q)\leftrightarrow(\neg P\lor Q)}
  {\FormulaRefAuto{P\to Q\eqvdash\neg P\lor Q}[thm]}
 \proofstepwidestar[]{(M,s\models p\to q)\leftrightarrow(P\to Q)}
  {\rLRS{5,\rLRS{4,1}}}
 \closeproofpartwideindR

 \proofpartwideindR[Bikonditional]{
 (M,s\models p\leftrightarrow q)\leftrightarrow(P\leftrightarrow Q)}
 {(M,s\models p\leftrightarrow q)\leftrightarrow(P\leftrightarrow Q)}
 [B48SatIffClause]
 \proofstepwidestar[]{\begin{aligned}
 (M,s\models p\leftrightarrow q)\leftrightarrow{}\\
 ((M,s\models p\to q)\land(M,s\models q\to p))
 \end{aligned}}
  {\FormulaRefAuto{B48LogicalAbbreviations}[def],
   \FormulaRefAuto{B48Satisfaction}[def]}
 \proofstepwidestar[]{(M,s\models p\to q)\leftrightarrow(P\to Q)}
  {\FormulaRefAuto{B48SatImpClause}[thm]}
 \proofstepwidestar[]{(M,s\models q\to p)\leftrightarrow(Q\to P)}
  {\FormulaRefAuto{B48SatImpClause}[thm]}
 \proofstepwidestar[]{(M,s\models p\leftrightarrow q)
                     \leftrightarrow((P\to Q)\land(Q\to P))}
  {\rLRS{2,\rLRS{3,1}}}
 \proofstepwidestar[]{(P\leftrightarrow Q)\leftrightarrow
                       ((P\to Q)\land(Q\to P))}
  {\FormulaRefAuto{P\leftrightarrow Q\eqvdash
                  (P\to Q)\land(Q\to P)}[thm]}
 \proofstepwidestar[]{(M,s\models p\leftrightarrow q)\leftrightarrow
                       (P\leftrightarrow Q)}
  {\rLRS{5,4}}
 \closeproofpartwideindR

 \proofpartwideindR[Allquantor]{(M,s\models\forall x\,p)\leftrightarrow
                \forall a\in D\;(M,s[x\mapsto a]\models p)}
 {(M,s\models\forall x\,p)\leftrightarrow
                \forall a\in D\;(M,s[x\mapsto a]\models p)}
 [B48SatAllClause]
 \proofstepwidestar[]{\begin{aligned}
 (M,s\models\forall x\,p)\leftrightarrow{}\\
 \neg\exists a\in D\;\neg(M,s[x\mapsto a]\models p)
 \end{aligned}}
  {\FormulaRefAuto{B48LogicalAbbreviations}[def],
   \FormulaRefAuto{B48Satisfaction}[def]}
 \proofstepwidestar[]{\begin{aligned}
 \forall a\;(a\in D\to(M,s[x\mapsto a]\models p))\leftrightarrow{}\\
 \neg\exists a\;(a\in D\land\neg(M,s[x\mapsto a]\models p))
 \end{aligned}}
  {\FormulaRefAuto{\forall x(P(x)\to Q(x))\eqvdash
                  \neg\exists x(P(x)\land\neg Q(x))}[thm]}
 \proofstepwidestar[]{(M,s\models\forall x\,p)\leftrightarrow
                \forall a\in D\;(M,s[x\mapsto a]\models p)}
  {\rLRS{2,1}}
 \closeproofpartwideindR
\end{tabproofsplitwide}

\FormulaDefDeltaK[Erfuellte Annahmen und Folgerung in einer Struktur]
{\begin{aligned}
 M,s\models\Gamma&\coloneq
       \forall p\in\Gamma\;(M,s\models p),\\
 \Gamma\models_M p&\coloneq
       \forall s\in A_D\;((M,s\models\Gamma)\to(M,s\models p)),\\
 x\mathrel{\#}\Gamma&\coloneq
       \forall q\in\Gamma\;(x\notin\operatorname{FV}(q)).
\end{aligned}}
{B48SemanticConsequence}{
 \DeltaRow{Struktur}{\mathsf{Str}(M)}
 \DeltaRow{Formeln und Annahmen}{p\in\mathcal F,\quad\Gamma\subseteq\mathcal F}
 \DeltaRow{Variable}{x\in\mathbb N}
}
Hier duerfen Annahmen freie Variablen haben. Die Quantifikation ueber
alle Belegungen ist Bestandteil der Folgerungsrelation.

\FormulaThmDeltaK[Anwendung einer semantischen Folgerung]
{\Gamma\models_M p,\ M,s\models\Gamma\ \vdash\ M,s\models p}
{B48ConsequenceInstance}{
 \DeltaRow{Struktur und Belegung}{\mathsf{Str}(M),\quad s\in A_D}
 \DeltaRow{Annahmen und Formel}{\Gamma\subseteq\mathcal F,\quad p\in\mathcal F}
}
\begin{tabproof}
 \proofstep{1}{\Gamma\models_M p}{\rA}
 \proofstep{2}{M,s\models\Gamma}{\rA}
 \proofstep{1}{\forall t\in A_D\;
    ((M,t\models\Gamma)\to(M,t\models p))}
  {\FormulaRefAuto{B48SemanticConsequence}[def]{1}}
 \proofstep{1}{(M,s\models\Gamma)\to(M,s\models p)}{\rUE{3}}
 \proofstep{1,2}{M,s\models p}{\rRE{4,2}}
\end{tabproof}

\FormulaThmDeltaK[Vereinigung erfuellter Annahmen]
{(M,s\models\Gamma\cup\Delta)\leftrightarrow
 ((M,s\models\Gamma)\land(M,s\models\Delta))}
{B48ContextUnion}{
 \DeltaRow{Struktur und Belegung}{\mathsf{Str}(M),\quad s\in A_D}
 \DeltaRow{Annahmenmengen}{\Gamma,\Delta\subseteq\mathcal F}
}
\begin{tabproof}
 \proofstep{1}{M,s\models\Gamma\cup\Delta}{\rA}
 \proofstep{2}{p\in\Gamma}{\rA}
 \proofstep{2}{p\in\Gamma\cup\Delta}
  {\FormulaRefAuto{x\in A\cup B\eqvdash x\in A\lor x\in B}[thm]{\rOIa{2}}}
 \proofstep{1,2}{M,s\models p}
  {\FormulaRefAuto{B48SemanticConsequence}[def]{1,3}}
 \proofstep{1}{M,s\models\Gamma}
  {\FormulaRefAuto{B48SemanticConsequence}[def]{\rUI{\rRI{2,4}}}}
 \proofstep{6}{p\in\Delta}{\rA}
 \proofstep{6}{p\in\Gamma\cup\Delta}
  {\FormulaRefAuto{x\in A\cup B\eqvdash x\in A\lor x\in B}[thm]{\rOIb{6}}}
 \proofstep{1,6}{M,s\models p}
  {\FormulaRefAuto{B48SemanticConsequence}[def]{1,7}}
 \proofstep{1}{M,s\models\Delta}
  {\FormulaRefAuto{B48SemanticConsequence}[def]{\rUI{\rRI{6,8}}}}
 \proofstep{1}{(M,s\models\Gamma)\land(M,s\models\Delta)}
  {\rAI{5,9}}
 \proofstep{}{(M,s\models\Gamma\cup\Delta)\to
        ((M,s\models\Gamma)\land(M,s\models\Delta))}
  {\rRI{1,10}}
 \proofstep{12}{(M,s\models\Gamma)\land(M,s\models\Delta)}{\rA}
 \proofstep{13}{p\in\Gamma\cup\Delta}{\rA}
 \proofstep{13}{p\in\Gamma\lor p\in\Delta}
  {\FormulaRefAuto{x\in A\cup B\eqvdash x\in A\lor x\in B}[thm]{13}}
 \proofstep{15}{p\in\Gamma}{\rA}
 \proofstep{12,15}{M,s\models p}
  {\FormulaRefAuto{B48SemanticConsequence}[def]{\rAEa{12},15}}
 \proofstep{17}{p\in\Delta}{\rA}
 \proofstep{12,17}{M,s\models p}
  {\FormulaRefAuto{B48SemanticConsequence}[def]{\rAEb{12},17}}
 \proofstep{12,13}{M,s\models p}{\rOE{14,15:16,17:18}}
 \proofstep{12}{M,s\models\Gamma\cup\Delta}
  {\FormulaRefAuto{B48SemanticConsequence}[def]{\rUI{\rRI{13,19}}}}
 \proofstep{}{((M,s\models\Gamma)\land(M,s\models\Delta))\to
                 (M,s\models\Gamma\cup\Delta)}
  {\rRI{12,20}}
 \proofstep{}{(M,s\models\Gamma\cup\Delta)\leftrightarrow
       ((M,s\models\Gamma)\land(M,s\models\Delta))}
  {\rLRI{11,21}}
\end{tabproof}

\FormulaThmDeltaK[Hinzunahme einer erfuellten Annahme]
{(M,s\models\Gamma\cup\{p\})\leftrightarrow
 ((M,s\models\Gamma)\land(M,s\models p))}
{B48ContextAssumption}{
 \DeltaRow{Struktur und Belegung}{\mathsf{Str}(M),\quad s\in A_D}
 \DeltaRow{Annahmen und Formel}{\Gamma\subseteq\mathcal F,\quad p\in\mathcal F}
}
\begin{tabproof}
 \proofstep{1}{M,s\models\{p\}}{\rA}
 \proofstep{1}{M,s\models p}
  {\FormulaRefAuto{B48SemanticConsequence}[def]{1}}
 \proofstep{}{(M,s\models\{p\})\to(M,s\models p)}{\rRI{1,2}}
 \proofstep{4}{M,s\models p}{\rA}
 \proofstep{5}{q\in\{p\}}{\rA}
 \proofstep{5}{q=p}{\FormulaRefAuto{x\in\{a\}\eqvdash x=a}[thm]{5}}
 \proofstep{4,5}{M,s\models q}
  {\rIE{\FormulaRefAuto{a=b\vdash b=a}[thm]{6},4}}
 \proofstep{4}{M,s\models\{p\}}
  {\FormulaRefAuto{B48SemanticConsequence}[def]{\rUI{\rRI{5,7}}}}
 \proofstep{}{(M,s\models p)\to(M,s\models\{p\})}{\rRI{4,8}}
 \proofstep{}{(M,s\models\{p\})\leftrightarrow(M,s\models p)}
  {\rLRI{3,9}}
 \proofstep{}{(M,s\models\Gamma\cup\{p\})\leftrightarrow
    ((M,s\models\Gamma)\land(M,s\models p))}
  {\rLRS{10,\FormulaRefAuto{B48ContextUnion}[thm]}}
\end{tabproof}

\FormulaThmDeltaK[Belegungswechsel in variablenfreien Annahmen]
{x\mathrel{\#}\Gamma,\ M,s\models\Gamma\ \vdash\
                         M,s[x\mapsto a]\models\Gamma}
{B48ContextFresh}{
 \DeltaRow{Struktur und Belegung}{\mathsf{Str}(M),\quad s\in A_D}
 \DeltaRow{Annahmen, Variable und Wert}{\Gamma\subseteq\mathcal F,\quad
                             x\in\mathbb N,\quad a\in D}
}
\begin{tabproof}
 \proofstep{1}{x\mathrel{\#}\Gamma}{\rA}
 \proofstep{2}{M,s\models\Gamma}{\rA}
 \proofstep{3}{p\in\Gamma}{\rA}
 \proofstep{1,3}{x\notin\operatorname{FV}(p)}
  {\FormulaRefAuto{B48SemanticConsequence}[def]{1,3}}
 \proofstep{2,3}{M,s\models p}
  {\FormulaRefAuto{B48SemanticConsequence}[def]{2,3}}
 \proofstep{1,3}{(M,s\models p)\leftrightarrow
                   (M,s[x\mapsto a]\models p)}
  {\FormulaRefAuto{B48FreshUpdate}[thm]{4}}
 \proofstep{1,2,3}{M,s[x\mapsto a]\models p}{\rLREa{6,5}}
 \proofstep{1,2}{M,s[x\mapsto a]\models\Gamma}
  {\FormulaRefAuto{B48SemanticConsequence}[def]{\rUI{\rRI{3,7}}}}
\end{tabproof}

\subsection{Einzeln nachgewiesene Regelfaelle}

In den folgenden Beweisen wird \(s\in A_D\) beliebig gewaehlt.
Die abschliessende All-Einfuehrung bindet diese Wahl. Die Formeln
\(p,q,r\) und die Annahmenmengen sind stets im zuvor angegebenen
Formelbereich.

\FormulaThmDeltaK[Semantischer Regelfall der ersten Konjunktionsbeseitigung]
{\Gamma\models_M p\land q\ \vdash\ \Gamma\models_M p}
{B48RuleAndLeft}{
 \DeltaRow{Struktur}{\mathsf{Str}(M)}
 \DeltaRow{Annahmen}{\Gamma,\Delta\subseteq\mathcal F}
 \DeltaRow{Formeln}{p,q,r\in\mathcal F}
}
\begin{tabproofsplitwide}
 \proofpartwide{Beweis}
 \proofstepwidestar[1]{\Gamma\models_M p\land q}
  {\rA}
 \proofstepwidestar[2]{M,s\models \Gamma}
  {\rA}
 \proofstepwidestar[1,2]{M,s\models p\land q}
  {\FormulaRefAuto{B48ConsequenceInstance}[thm]{1,2}}
 \proofstepwidestar[1,2]{M,s\models p}
  {\FormulaRefAuto{B48SatAndLeft}[thm]{3}}
 \proofstepwidestar[1]{(M,s\models \Gamma)\to(M,s\models p)}
  {\rRI{2,4}}
 \proofstepwidestar[1]{\Gamma\models_M p}
  {\FormulaRefAuto{B48SemanticConsequence}[def]{\rUI{5}}}
 \closeproofpartwide
\end{tabproofsplitwide}

\FormulaThmDeltaK[Semantischer Regelfall der zweiten Konjunktionsbeseitigung]
{\Gamma\models_M p\land q\ \vdash\ \Gamma\models_M q}
{B48RuleAndRight}{
 \DeltaRow{Struktur}{\mathsf{Str}(M)}
 \DeltaRow{Annahmen}{\Gamma,\Delta\subseteq\mathcal F}
 \DeltaRow{Formeln}{p,q,r\in\mathcal F}
}
\begin{tabproofsplitwide}
 \proofpartwide{Beweis}
 \proofstepwidestar[1]{\Gamma\models_M p\land q}
  {\rA}
 \proofstepwidestar[2]{M,s\models \Gamma}
  {\rA}
 \proofstepwidestar[1,2]{M,s\models p\land q}
  {\FormulaRefAuto{B48ConsequenceInstance}[thm]{1,2}}
 \proofstepwidestar[1,2]{M,s\models q}
  {\FormulaRefAuto{B48SatAndRight}[thm]{3}}
 \proofstepwidestar[1]{(M,s\models \Gamma)\to(M,s\models q)}
  {\rRI{2,4}}
 \proofstepwidestar[1]{\Gamma\models_M q}
  {\FormulaRefAuto{B48SemanticConsequence}[def]{\rUI{5}}}
 \closeproofpartwide
\end{tabproofsplitwide}

\FormulaThmDeltaK[Semantische Disjunktionseinfuehrung links]
{\Gamma\models_M p\ \vdash\ \Gamma\models_M p\lor q}
{B48RuleOrLeft}{
 \DeltaRow{Struktur}{\mathsf{Str}(M)}
 \DeltaRow{Annahmen}{\Gamma,\Delta\subseteq\mathcal F}
 \DeltaRow{Formeln}{p,q,r\in\mathcal F}
}
\begin{tabproofsplitwide}
 \proofpartwide{Beweis}
 \proofstepwidestar[1]{\Gamma\models_M p}
  {\rA}
 \proofstepwidestar[2]{M,s\models \Gamma}
  {\rA}
 \proofstepwidestar[1,2]{M,s\models p}
  {\FormulaRefAuto{B48ConsequenceInstance}[thm]{1,2}}
 \proofstepwidestar[1,2]{(M,s\models p)\lor(M,s\models q)}
  {\rOIa{3}}
 \proofstepwidestar[1,2]{M,s\models p\lor q}
  {\rLREb{\FormulaRefAuto{B48SatOrClause}[thm],4}}
 \proofstepwidestar[1]{(M,s\models \Gamma)\to(M,s\models p\lor q)}
  {\rRI{2,5}}
 \proofstepwidestar[1]{\Gamma\models_M p\lor q}
  {\FormulaRefAuto{B48SemanticConsequence}[def]{\rUI{6}}}
 \closeproofpartwide
\end{tabproofsplitwide}

\FormulaThmDeltaK[Semantische Disjunktionseinfuehrung rechts]
{\Gamma\models_M q\ \vdash\ \Gamma\models_M p\lor q}
{B48RuleOrRight}{
 \DeltaRow{Struktur}{\mathsf{Str}(M)}
 \DeltaRow{Annahmen}{\Gamma,\Delta\subseteq\mathcal F}
 \DeltaRow{Formeln}{p,q,r\in\mathcal F}
}
\begin{tabproofsplitwide}
 \proofpartwide{Beweis}
 \proofstepwidestar[1]{\Gamma\models_M q}
  {\rA}
 \proofstepwidestar[2]{M,s\models \Gamma}
  {\rA}
 \proofstepwidestar[1,2]{M,s\models q}
  {\FormulaRefAuto{B48ConsequenceInstance}[thm]{1,2}}
 \proofstepwidestar[1,2]{(M,s\models p)\lor(M,s\models q)}
  {\rOIb{3}}
 \proofstepwidestar[1,2]{M,s\models p\lor q}
  {\rLREb{\FormulaRefAuto{B48SatOrClause}[thm],4}}
 \proofstepwidestar[1]{(M,s\models \Gamma)\to(M,s\models p\lor q)}
  {\rRI{2,5}}
 \proofstepwidestar[1]{\Gamma\models_M p\lor q}
  {\FormulaRefAuto{B48SemanticConsequence}[def]{\rUI{6}}}
 \closeproofpartwide
\end{tabproofsplitwide}

\FormulaThmDeltaK[Semantische Bikonditionalbeseitigung vorwaerts]
{\Gamma\models_M p\leftrightarrow q\ \vdash\ \Gamma\models_M p\to q}
{B48RuleIffLeft}{
 \DeltaRow{Struktur}{\mathsf{Str}(M)}
 \DeltaRow{Annahmen}{\Gamma,\Delta\subseteq\mathcal F}
 \DeltaRow{Formeln}{p,q,r\in\mathcal F}
}
\begin{tabproofsplitwide}
 \proofpartwide{Beweis}
 \proofstepwidestar[1]{\Gamma\models_M p\leftrightarrow q}
  {\rA}
 \proofstepwidestar[2]{M,s\models \Gamma}
  {\rA}
 \proofstepwidestar[1,2]{M,s\models p\leftrightarrow q}
  {\FormulaRefAuto{B48ConsequenceInstance}[thm]{1,2}}
 \proofstepwidestar[1,2]{(M,s\models p)\leftrightarrow(M,s\models q)}
  {\rLREa{\FormulaRefAuto{B48SatIffClause}[thm],3}}
 \proofstepwidestar[1,2]{(M,s\models p)\to(M,s\models q)}
  {\rLREa{4}}
 \proofstepwidestar[1,2]{M,s\models p\to q}
  {\rLREb{\FormulaRefAuto{B48SatImpClause}[thm],5}}
 \proofstepwidestar[1]{(M,s\models \Gamma)\to(M,s\models p\to q)}
  {\rRI{2,6}}
 \proofstepwidestar[1]{\Gamma\models_M p\to q}
  {\FormulaRefAuto{B48SemanticConsequence}[def]{\rUI{7}}}
 \closeproofpartwide
\end{tabproofsplitwide}

\FormulaThmDeltaK[Semantische Bikonditionalbeseitigung rueckwaerts]
{\Gamma\models_M p\leftrightarrow q\ \vdash\ \Gamma\models_M q\to p}
{B48RuleIffRight}{
 \DeltaRow{Struktur}{\mathsf{Str}(M)}
 \DeltaRow{Annahmen}{\Gamma,\Delta\subseteq\mathcal F}
 \DeltaRow{Formeln}{p,q,r\in\mathcal F}
}
\begin{tabproofsplitwide}
 \proofpartwide{Beweis}
 \proofstepwidestar[1]{\Gamma\models_M p\leftrightarrow q}
  {\rA}
 \proofstepwidestar[2]{M,s\models \Gamma}
  {\rA}
 \proofstepwidestar[1,2]{M,s\models p\leftrightarrow q}
  {\FormulaRefAuto{B48ConsequenceInstance}[thm]{1,2}}
 \proofstepwidestar[1,2]{(M,s\models p)\leftrightarrow(M,s\models q)}
  {\rLREa{\FormulaRefAuto{B48SatIffClause}[thm],3}}
 \proofstepwidestar[1,2]{(M,s\models q)\to(M,s\models p)}
  {\rLREb{4}}
 \proofstepwidestar[1,2]{M,s\models q\to p}
  {\rLREb{\FormulaRefAuto{B48SatImpClause}[thm],5}}
 \proofstepwidestar[1]{(M,s\models \Gamma)\to(M,s\models q\to p)}
  {\rRI{2,6}}
 \proofstepwidestar[1]{\Gamma\models_M q\to p}
  {\FormulaRefAuto{B48SemanticConsequence}[def]{\rUI{7}}}
 \closeproofpartwide
\end{tabproofsplitwide}

\FormulaThmDeltaK[Semantische Konjunktionseinfuehrung fuer Annahmenmengen]
{\begin{gathered}\Gamma\models_M p,\quad\Delta\models_M q\\
 \vdash\ \Gamma\cup\Delta\models_M p\land q\end{gathered}}
{B48RuleAndIntro}{
 \DeltaRow{Struktur}{\mathsf{Str}(M)}
 \DeltaRow{Annahmen}{\Gamma,\Delta\subseteq\mathcal F}
 \DeltaRow{Formeln}{p,q,r\in\mathcal F}
}
\begin{tabproofsplitwide}
 \proofpartwide{Beweis}
 \proofstepwidestar[1]{\Gamma\models_M p}
  {\rA}
 \proofstepwidestar[2]{\Delta\models_M q}
  {\rA}
 \proofstepwidestar[3]{M,s\models \Gamma\cup\Delta}
  {\rA}
 \proofstepwidestar[3]{(M,s\models \Gamma)\land(M,s\models \Delta)}
  {\rLREa{\FormulaRefAuto{B48ContextUnion}[thm],3}}
 \proofstepwidestar[3]{M,s\models \Gamma}
  {\rAEa{4}}
 \proofstepwidestar[3]{M,s\models \Delta}
  {\rAEb{4}}
 \proofstepwidestar[1,3]{M,s\models p}
  {\FormulaRefAuto{B48ConsequenceInstance}[thm]{1,5}}
 \proofstepwidestar[2,3]{M,s\models q}
  {\FormulaRefAuto{B48ConsequenceInstance}[thm]{2,6}}
 \proofstepwidestar[1,2,3]{M,s\models p\land q}
  {\FormulaRefAuto{B48SatAndIntro}[thm]{7,8}}
 \proofstepwidestar[1,2]{(M,s\models \Gamma\cup\Delta)\to(M,s\models p\land q)}
  {\rRI{3,9}}
 \proofstepwidestar[1,2]{\Gamma\cup\Delta\models_M p\land q}
  {\FormulaRefAuto{B48SemanticConsequence}[def]{\rUI{10}}}
 \closeproofpartwide
\end{tabproofsplitwide}

\FormulaThmDeltaK[Semantische Implikationsbeseitigung]
{\begin{gathered}\Gamma\models_M p\to q,\quad\Delta\models_M p\\
 \vdash\ \Gamma\cup\Delta\models_M q\end{gathered}}
{B48RuleImpElim}{
 \DeltaRow{Struktur}{\mathsf{Str}(M)}
 \DeltaRow{Annahmen}{\Gamma,\Delta\subseteq\mathcal F}
 \DeltaRow{Formeln}{p,q,r\in\mathcal F}
}
\begin{tabproofsplitwide}
 \proofpartwide{Beweis}
 \proofstepwidestar[1]{\Gamma\models_M p\to q}
  {\rA}
 \proofstepwidestar[2]{\Delta\models_M p}
  {\rA}
 \proofstepwidestar[3]{M,s\models \Gamma\cup\Delta}
  {\rA}
 \proofstepwidestar[3]{(M,s\models \Gamma)\land(M,s\models \Delta)}
  {\rLREa{\FormulaRefAuto{B48ContextUnion}[thm],3}}
 \proofstepwidestar[3]{M,s\models \Gamma}
  {\rAEa{4}}
 \proofstepwidestar[3]{M,s\models \Delta}
  {\rAEb{4}}
 \proofstepwidestar[1,3]{M,s\models p\to q}
  {\FormulaRefAuto{B48ConsequenceInstance}[thm]{1,5}}
 \proofstepwidestar[2,3]{M,s\models p}
  {\FormulaRefAuto{B48ConsequenceInstance}[thm]{2,6}}
 \proofstepwidestar[1,3]{(M,s\models p)\to(M,s\models q)}
  {\rLREa{\FormulaRefAuto{B48SatImpClause}[thm],7}}
 \proofstepwidestar[1,2,3]{M,s\models q}
  {\rRE{9,8}}
 \proofstepwidestar[1,2]{(M,s\models \Gamma\cup\Delta)\to(M,s\models q)}
  {\rRI{3,10}}
 \proofstepwidestar[1,2]{\Gamma\cup\Delta\models_M q}
  {\FormulaRefAuto{B48SemanticConsequence}[def]{\rUI{11}}}
 \closeproofpartwide
\end{tabproofsplitwide}

\FormulaThmDeltaK[Semantische Bikonditionaleinfuehrung]
{\begin{gathered}\Gamma\models_M p\to q,\quad\Delta\models_M q\to p\\
 \vdash\ \Gamma\cup\Delta\models_M p\leftrightarrow q\end{gathered}}
{B48RuleIffIntro}{
 \DeltaRow{Struktur}{\mathsf{Str}(M)}
 \DeltaRow{Annahmen}{\Gamma,\Delta\subseteq\mathcal F}
 \DeltaRow{Formeln}{p,q,r\in\mathcal F}
}
\begin{tabproofsplitwide}
 \proofpartwide{Beweis}
 \proofstepwidestar[1]{\Gamma\models_M p\to q}
  {\rA}
 \proofstepwidestar[2]{\Delta\models_M q\to p}
  {\rA}
 \proofstepwidestar[3]{M,s\models \Gamma\cup\Delta}
  {\rA}
 \proofstepwidestar[3]{(M,s\models \Gamma)\land(M,s\models \Delta)}
  {\rLREa{\FormulaRefAuto{B48ContextUnion}[thm],3}}
 \proofstepwidestar[3]{M,s\models \Gamma}
  {\rAEa{4}}
 \proofstepwidestar[3]{M,s\models \Delta}
  {\rAEb{4}}
 \proofstepwidestar[1,3]{M,s\models p\to q}
  {\FormulaRefAuto{B48ConsequenceInstance}[thm]{1,5}}
 \proofstepwidestar[2,3]{M,s\models q\to p}
  {\FormulaRefAuto{B48ConsequenceInstance}[thm]{2,6}}
 \proofstepwidestar[1,3]{(M,s\models p)\to(M,s\models q)}
  {\rLREa{\FormulaRefAuto{B48SatImpClause}[thm],7}}
 \proofstepwidestar[2,3]{(M,s\models q)\to(M,s\models p)}
  {\rLREa{\FormulaRefAuto{B48SatImpClause}[thm],8}}
 \proofstepwidestar[1,2,3]{(M,s\models p)\leftrightarrow(M,s\models q)}
  {\rLRI{9,10}}
 \proofstepwidestar[1,2,3]{M,s\models p\leftrightarrow q}
  {\rLREb{\FormulaRefAuto{B48SatIffClause}[thm],11}}
 \proofstepwidestar[1,2]{(M,s\models \Gamma\cup\Delta)\to(M,s\models p\leftrightarrow q)}
  {\rRI{3,12}}
 \proofstepwidestar[1,2]{\Gamma\cup\Delta\models_M p\leftrightarrow q}
  {\FormulaRefAuto{B48SemanticConsequence}[def]{\rUI{13}}}
 \closeproofpartwide
\end{tabproofsplitwide}

\FormulaThmDeltaK[Semantische Widerspruchseinfuehrung]
{\begin{gathered}\Gamma\models_M p,\quad\Delta\models_M \neg p\\
 \vdash\ \Gamma\cup\Delta\models_M \bot\end{gathered}}
{B48RuleBottom}{
 \DeltaRow{Struktur}{\mathsf{Str}(M)}
 \DeltaRow{Annahmen}{\Gamma,\Delta\subseteq\mathcal F}
 \DeltaRow{Formeln}{p,q,r\in\mathcal F}
}
\begin{tabproofsplitwide}
 \proofpartwide{Beweis}
 \proofstepwidestar[1]{\Gamma\models_M p}
  {\rA}
 \proofstepwidestar[2]{\Delta\models_M \neg p}
  {\rA}
 \proofstepwidestar[3]{M,s\models \Gamma\cup\Delta}
  {\rA}
 \proofstepwidestar[3]{(M,s\models \Gamma)\land(M,s\models \Delta)}
  {\rLREa{\FormulaRefAuto{B48ContextUnion}[thm],3}}
 \proofstepwidestar[3]{M,s\models \Gamma}
  {\rAEa{4}}
 \proofstepwidestar[3]{M,s\models \Delta}
  {\rAEb{4}}
 \proofstepwidestar[1,3]{M,s\models p}
  {\FormulaRefAuto{B48ConsequenceInstance}[thm]{1,5}}
 \proofstepwidestar[2,3]{M,s\models \neg p}
  {\FormulaRefAuto{B48ConsequenceInstance}[thm]{2,6}}
 \proofstepwidestar[1,2,3]{\bot}
  {\FormulaRefAuto{B48SatContradiction}[thm]{7,8}}
 \proofstepwidestar[1,2,3]{M,s\models \bot}
  {\FormulaRefAuto{B48Satisfaction}[def]{9}}
 \proofstepwidestar[1,2]{(M,s\models \Gamma\cup\Delta)\to(M,s\models \bot)}
  {\rRI{3,10}}
 \proofstepwidestar[1,2]{\Gamma\cup\Delta\models_M \bot}
  {\FormulaRefAuto{B48SemanticConsequence}[def]{\rUI{11}}}
 \closeproofpartwide
\end{tabproofsplitwide}

\FormulaThmDeltaK[Semantische Implikationseinfuehrung]
{\Gamma\cup\{p\}\models_M q\ \vdash\ \Gamma\models_M p\to q}
{B48RuleImpIntro}{
 \DeltaRow{Struktur}{\mathsf{Str}(M)}
 \DeltaRow{Annahmen}{\Gamma,\Delta,\Theta\subseteq\mathcal F}
 \DeltaRow{Formeln und Variablen}{p,q,r\in\mathcal F,\quad x,y,z\in\mathbb N}
}
\begin{tabproofsplitwide}
 \proofpartwide{Beweis}
 \proofstepwidestar[1]{\Gamma\cup\{p\}\models_M q}
  {\rA}
 \proofstepwidestar[2]{M,s\models \Gamma}
  {\rA}
 \proofstepwidestar[3]{M,s\models p}
  {\rA}
 \proofstepwidestar[2,3]{(M,s\models \Gamma)\land(M,s\models p)}
  {\rAI{2,3}}
 \proofstepwidestar[2,3]{M,s\models \Gamma\cup\{p\}}
  {\rLREb{\FormulaRefAuto{B48ContextAssumption}[thm],4}}
 \proofstepwidestar[1,2,3]{M,s\models q}
  {\FormulaRefAuto{B48ConsequenceInstance}[thm]{1,5}}
 \proofstepwidestar[1,2]{(M,s\models p)\to(M,s\models q)}
  {\rRI{3,6}}
 \proofstepwidestar[1,2]{M,s\models p\to q}
  {\rLREb{\FormulaRefAuto{B48SatImpClause}[thm],7}}
 \proofstepwidestar[1]{(M,s\models \Gamma)\to(M,s\models p\to q)}
  {\rRI{2,8}}
 \proofstepwidestar[1]{\Gamma\models_M p\to q}
  {\FormulaRefAuto{B48SemanticConsequence}[def]{\rUI{9}}}
 \closeproofpartwide
\end{tabproofsplitwide}

\FormulaThmDeltaK[Semantische Negationseinfuehrung mit Entladung]
{\Gamma\cup\{p\}\models_M \bot\ \vdash\ \Gamma\models_M \neg p}
{B48RuleNegIntro}{
 \DeltaRow{Struktur}{\mathsf{Str}(M)}
 \DeltaRow{Annahmen}{\Gamma,\Delta,\Theta\subseteq\mathcal F}
 \DeltaRow{Formeln und Variablen}{p,q,r\in\mathcal F,\quad x,y,z\in\mathbb N}
}
\begin{tabproofsplitwide}
 \proofpartwide{Beweis}
 \proofstepwidestar[1]{\Gamma\cup\{p\}\models_M \bot}
  {\rA}
 \proofstepwidestar[2]{M,s\models \Gamma}
  {\rA}
 \proofstepwidestar[3]{M,s\models p}
  {\rA}
 \proofstepwidestar[2,3]{(M,s\models \Gamma)\land(M,s\models p)}
  {\rAI{2,3}}
 \proofstepwidestar[2,3]{M,s\models \Gamma\cup\{p\}}
  {\rLREb{\FormulaRefAuto{B48ContextAssumption}[thm],4}}
 \proofstepwidestar[1,2,3]{M,s\models \bot}
  {\FormulaRefAuto{B48ConsequenceInstance}[thm]{1,5}}
 \proofstepwidestar[1,2,3]{\bot}
  {\FormulaRefAuto{B48Satisfaction}[def]{6}}
 \proofstepwidestar[1,2]{\neg(M,s\models p)}
  {\rCI{3,7}}
 \proofstepwidestar[1,2]{M,s\models \neg p}
  {\FormulaRefAuto{B48Satisfaction}[def]{8}}
 \proofstepwidestar[1]{(M,s\models \Gamma)\to(M,s\models \neg p)}
  {\rRI{2,9}}
 \proofstepwidestar[1]{\Gamma\models_M \neg p}
  {\FormulaRefAuto{B48SemanticConsequence}[def]{\rUI{10}}}
 \closeproofpartwide
\end{tabproofsplitwide}

\FormulaThmDeltaK[Semantischer klassischer Regelfall mit Entladung]
{\Gamma\cup\{\neg p\}\models_M \bot\ \vdash\ \Gamma\models_M p}
{B48RuleClassical}{
 \DeltaRow{Struktur}{\mathsf{Str}(M)}
 \DeltaRow{Annahmen}{\Gamma,\Delta,\Theta\subseteq\mathcal F}
 \DeltaRow{Formeln und Variablen}{p,q,r\in\mathcal F,\quad x,y,z\in\mathbb N}
}
\begin{tabproofsplitwide}
 \proofpartwide{Beweis}
 \proofstepwidestar[1]{\Gamma\cup\{\neg p\}\models_M \bot}
  {\rA}
 \proofstepwidestar[2]{M,s\models \Gamma}
  {\rA}
 \proofstepwidestar[3]{\neg(M,s\models p)}
  {\rA}
 \proofstepwidestar[3]{M,s\models \neg p}
  {\FormulaRefAuto{B48Satisfaction}[def]{3}}
 \proofstepwidestar[2,3]{(M,s\models \Gamma)\land(M,s\models \neg p)}
  {\rAI{2,4}}
 \proofstepwidestar[2,3]{M,s\models \Gamma\cup\{\neg p\}}
  {\rLREb{\FormulaRefAuto{B48ContextAssumption}[thm],5}}
 \proofstepwidestar[1,2,3]{M,s\models \bot}
  {\FormulaRefAuto{B48ConsequenceInstance}[thm]{1,6}}
 \proofstepwidestar[1,2,3]{\bot}
  {\FormulaRefAuto{B48Satisfaction}[def]{7}}
 \proofstepwidestar[1,2]{M,s\models p}
  {\rCE{3,8}}
 \proofstepwidestar[1]{(M,s\models \Gamma)\to(M,s\models p)}
  {\rRI{2,9}}
 \proofstepwidestar[1]{\Gamma\models_M p}
  {\FormulaRefAuto{B48SemanticConsequence}[def]{\rUI{10}}}
 \closeproofpartwide
\end{tabproofsplitwide}

\FormulaThmDeltaK[Semantische Disjunktionsbeseitigung mit beiden Entladungen]
{\begin{gathered}
 \Gamma\models_M p\lor q,\quad
 \Delta\cup\{p\}\models_M r,\quad\Theta\cup\{q\}\models_M r\\
 \vdash\ (\Gamma\cup\Delta)\cup\Theta\models_M r
 \end{gathered}}
{B48RuleOrElim}{
 \DeltaRow{Struktur}{\mathsf{Str}(M)}
 \DeltaRow{Annahmen}{\Gamma,\Delta,\Theta\subseteq\mathcal F}
 \DeltaRow{Formeln und Variablen}{p,q,r\in\mathcal F,\quad x,y,z\in\mathbb N}
}
\begin{tabproofsplitwide}
 \proofpartwide{Beweis}
 \proofstepwidestar[1]{\Gamma\models_M p\lor q}
  {\rA}
 \proofstepwidestar[2]{\Delta\cup\{p\}\models_M r}
  {\rA}
 \proofstepwidestar[3]{\Theta\cup\{q\}\models_M r}
  {\rA}
 \proofstepwidestar[4]{M,s\models (\Gamma\cup\Delta)\cup\Theta}
  {\rA}
 \proofstepwidestar[4]{(M,s\models \Gamma\cup\Delta)\land(M,s\models \Theta)}
  {\rLREa{\FormulaRefAuto{B48ContextUnion}[thm],4}}
 \proofstepwidestar[4]{(M,s\models \Gamma)\land(M,s\models \Delta)}
  {\rLREa{\FormulaRefAuto{B48ContextUnion}[thm],\rAEa{5}}}
 \proofstepwidestar[1,4]{M,s\models p\lor q}
  {\FormulaRefAuto{B48ConsequenceInstance}[thm]{1,\rAEa{6}}}
 \proofstepwidestar[1,4]{(M,s\models p)\lor(M,s\models q)}
  {\rLREa{\FormulaRefAuto{B48SatOrClause}[thm],7}}
 \proofstepwidestar[9]{M,s\models p}
  {\rA}
 \proofstepwidestar[4,9]{(M,s\models \Delta)\land(M,s\models p)}
  {\rAI{\rAEb{6},9}}
 \proofstepwidestar[4,9]{M,s\models \Delta\cup\{p\}}
  {\rLREb{\FormulaRefAuto{B48ContextAssumption}[thm],10}}
 \proofstepwidestar[2,4,9]{M,s\models r}
  {\FormulaRefAuto{B48ConsequenceInstance}[thm]{2,11}}
 \proofstepwidestar[13]{M,s\models q}
  {\rA}
 \proofstepwidestar[4,13]{(M,s\models \Theta)\land(M,s\models q)}
  {\rAI{\rAEb{5},13}}
 \proofstepwidestar[4,13]{M,s\models \Theta\cup\{q\}}
  {\rLREb{\FormulaRefAuto{B48ContextAssumption}[thm],14}}
 \proofstepwidestar[3,4,13]{M,s\models r}
  {\FormulaRefAuto{B48ConsequenceInstance}[thm]{3,15}}
 \proofstepwidestar[1,2,3,4]{M,s\models r}
  {\rOE{8,9:12,13:16}}
 \proofstepwidestar[1,2,3]{(M,s\models (\Gamma\cup\Delta)\cup\Theta)\to(M,s\models r)}
  {\rRI{4,17}}
 \proofstepwidestar[1,2,3]{(\Gamma\cup\Delta)\cup\Theta\models_M r}
  {\FormulaRefAuto{B48SemanticConsequence}[def]{\rUI{18}}}
 \closeproofpartwide
\end{tabproofsplitwide}

\FormulaThmDeltaK[Semantische All-Einfuehrung]
{\Gamma\models_M p,\quad x\mathrel{\#}\Gamma\
   \vdash\ \Gamma\models_M\forall x\,p}
{B48RuleAllIntro}{
 \DeltaRow{Struktur}{\mathsf{Str}(M)}
 \DeltaRow{Annahmen}{\Gamma,\Delta,\Theta\subseteq\mathcal F}
 \DeltaRow{Formeln und Variablen}{p,q,r\in\mathcal F,\quad x,y,z\in\mathbb N}
}
\begin{tabproofsplitwide}
 \proofpartwide{Beweis}
 \proofstepwidestar[1]{\Gamma\models_M p}
  {\rA}
 \proofstepwidestar[2]{x\mathrel{\#}\Gamma}
  {\rA}
 \proofstepwidestar[3]{M,s\models \Gamma}
  {\rA}
 \proofstepwidestar[4]{a\in D}
  {\rA}
 \proofstepwidestar[4]{s[x\mapsto a]\in A_D}
  {\FormulaRefAuto{B48AssignmentUpdate}[thm]{4}}
 \proofstepwidestar[2,3,4]{M,s[x\mapsto a]\models\Gamma}
  {\FormulaRefAuto{B48ContextFresh}[thm]{2,3,4}}
 \proofstepwidestar[1,2,3,4]{M,s[x\mapsto a]\models p}
  {\FormulaRefAuto{B48ConsequenceInstance}[thm]{1,6,5}}
 \proofstepwidestar[1,2,3]{\forall a\in D\;(M,s[x\mapsto a]\models p)}
  {\rUI{\rRI{4,7}}}
 \proofstepwidestar[1,2,3]{M,s\models \forall x\,p}
  {\rLREb{\FormulaRefAuto{B48SatAllClause}[thm],8}}
 \proofstepwidestar[1,2]{(M,s\models \Gamma)\to(M,s\models \forall x\,p)}
  {\rRI{3,9}}
 \proofstepwidestar[1,2]{\Gamma\models_M \forall x\,p}
  {\FormulaRefAuto{B48SemanticConsequence}[def]{\rUI{10}}}
 \closeproofpartwide
\end{tabproofsplitwide}

\FormulaThmDeltaK[Semantische All-Beseitigung mit Substitution]
{\Gamma\models_M \forall x\,p,\quad\mathsf{Fr}(p,x,y)\
 \vdash\ \Gamma\models_M p[y/x]}
{B48RuleAllElim}{
 \DeltaRow{Struktur}{\mathsf{Str}(M)}
 \DeltaRow{Annahmen}{\Gamma,\Delta,\Theta\subseteq\mathcal F}
 \DeltaRow{Formeln und Variablen}{p,q,r\in\mathcal F,\quad x,y,z\in\mathbb N}
}
\begin{tabproofsplitwide}
 \proofpartwide{Beweis}
 \proofstepwidestar[1]{\Gamma\models_M \forall x\,p}
  {\rA}
 \proofstepwidestar[2]{\mathsf{Fr}(p,x,y)}
  {\rA}
 \proofstepwidestar[3]{M,s\models \Gamma}
  {\rA}
 \proofstepwidestar[1,3]{M,s\models \forall x\,p}
  {\FormulaRefAuto{B48ConsequenceInstance}[thm]{1,3}}
 \proofstepwidestar[]{s\colon\mathbb N\to D}
  {\FormulaRefAuto{B48AssignmentCriterion}[thm]}
 \proofstepwidestar[]{s(y)\in D}
  {\FormulaRefAuto{F\colon A\to B,\ x\in A\vdash F(x)\in B}[thm]{5}}
 \proofstepwidestar[2]{(M,s\models p[y/x])\leftrightarrow(M,s[x\mapsto s(y)]\models p)}
  {\FormulaRefAuto{B48SubstitutionLemma}[thm]{2}}
 \proofstepwidestar[1,3]{\forall a\in D\;(M,s[x\mapsto a]\models p)}
  {\rLREa{\FormulaRefAuto{B48SatAllClause}[thm],4}}
 \proofstepwidestar[1,3]{M,s[x\mapsto s(y)]\models p}
  {\rRE{\rUE{8},6}}
 \proofstepwidestar[1,2,3]{M,s\models p[y/x]}
  {\rLREb{7,9}}
 \proofstepwidestar[1,2]{(M,s\models \Gamma)\to(M,s\models p[y/x])}
  {\rRI{3,10}}
 \proofstepwidestar[1,2]{\Gamma\models_M p[y/x]}
  {\FormulaRefAuto{B48SemanticConsequence}[def]{\rUI{11}}}
 \closeproofpartwide
\end{tabproofsplitwide}

\FormulaThmDeltaK[Semantische Existenz-Einfuehrung mit Substitution]
{\Gamma\models_M p[y/x],\quad\mathsf{Fr}(p,x,y)\
 \vdash\ \Gamma\models_M \exists x\,p}
{B48RuleExistsIntro}{
 \DeltaRow{Struktur}{\mathsf{Str}(M)}
 \DeltaRow{Annahmen}{\Gamma,\Delta,\Theta\subseteq\mathcal F}
 \DeltaRow{Formeln und Variablen}{p,q,r\in\mathcal F,\quad x,y,z\in\mathbb N}
}
\begin{tabproofsplitwide}
 \proofpartwide{Beweis}
 \proofstepwidestar[1]{\Gamma\models_M p[y/x]}
  {\rA}
 \proofstepwidestar[2]{\mathsf{Fr}(p,x,y)}
  {\rA}
 \proofstepwidestar[3]{M,s\models \Gamma}
  {\rA}
 \proofstepwidestar[1,3]{M,s\models p[y/x]}
  {\FormulaRefAuto{B48ConsequenceInstance}[thm]{1,3}}
 \proofstepwidestar[]{s\colon\mathbb N\to D}
  {\FormulaRefAuto{B48AssignmentCriterion}[thm]}
 \proofstepwidestar[]{s(y)\in D}
  {\FormulaRefAuto{F\colon A\to B,\ x\in A\vdash F(x)\in B}[thm]{5}}
 \proofstepwidestar[2]{(M,s\models p[y/x])\leftrightarrow(M,s[x\mapsto s(y)]\models p)}
  {\FormulaRefAuto{B48SubstitutionLemma}[thm]{2}}
 \proofstepwidestar[1,2,3]{M,s[x\mapsto s(y)]\models p}
  {\rLREa{7,4}}
 \proofstepwidestar[1,2,3]{M,s\models \exists x\,p}
  {\FormulaRefAuto{B48SatExistsIntro}[thm]{6,8}}
 \proofstepwidestar[1,2]{(M,s\models \Gamma)\to(M,s\models \exists x\,p)}
  {\rRI{3,9}}
 \proofstepwidestar[1,2]{\Gamma\models_M \exists x\,p}
  {\FormulaRefAuto{B48SemanticConsequence}[def]{\rUI{10}}}
 \closeproofpartwide
\end{tabproofsplitwide}

\FormulaThmDeltaK[Semantische Existenz-Beseitigung mit Eigenvariable]
{\begin{gathered}
 \Gamma\models_M\exists x\,p,\quad\Delta\cup\{p\}\models_M q,\\
 x\mathrel{\#}\Delta,\quad x\notin\operatorname{FV}(q)\\
 \vdash\ \Gamma\cup\Delta\models_M q
 \end{gathered}}
{B48RuleExistsElim}{
 \DeltaRow{Struktur}{\mathsf{Str}(M)}
 \DeltaRow{Annahmen}{\Gamma,\Delta,\Theta\subseteq\mathcal F}
 \DeltaRow{Formeln und Variablen}{p,q,r\in\mathcal F,\quad x,y,z\in\mathbb N}
}
Der Zeuge \(a\) ist frisch fuer die ausserhalb seines Teilbeweises offenen Annahmen.
\begin{tabproofsplitwide}
 \proofpartwide{Beweis}
 \proofstepwidestar[1]{\Gamma\models_M\exists x\,p}
  {\rA}
 \proofstepwidestar[2]{\Delta\cup\{p\}\models_M q}
  {\rA}
 \proofstepwidestar[3]{x\mathrel{\#}\Delta}
  {\rA}
 \proofstepwidestar[4]{x\notin\operatorname{FV}(q)}
  {\rA}
 \proofstepwidestar[5]{M,s\models \Gamma\cup\Delta}
  {\rA}
 \proofstepwidestar[5]{(M,s\models \Gamma)\land(M,s\models \Delta)}
  {\rLREa{\FormulaRefAuto{B48ContextUnion}[thm],5}}
 \proofstepwidestar[1,5]{M,s\models \exists x\,p}
  {\FormulaRefAuto{B48ConsequenceInstance}[thm]{1,\rAEa{6}}}
 \proofstepwidestar[1,5]{\exists a\;(a\in D\land(M,s[x\mapsto a]\models p))}
  {\FormulaRefAuto{B48Satisfaction}[def]{7}}
 \proofstepwidestar[9]{a\in D\land(M,s[x\mapsto a]\models p)}
  {\rA}
 \proofstepwidestar[9]{a\in D}
  {\rAEa{9}}
 \proofstepwidestar[9]{M,s[x\mapsto a]\models p}
  {\rAEb{9}}
 \proofstepwidestar[9]{s[x\mapsto a]\in A_D}
  {\FormulaRefAuto{B48AssignmentUpdate}[thm]{10}}
 \proofstepwidestar[3,5,9]{M,s[x\mapsto a]\models\Delta}
  {\FormulaRefAuto{B48ContextFresh}[thm]{3,\rAEb{6},10}}
 \proofstepwidestar[3,5,9]{M,s[x\mapsto a]\models\Delta\cup\{p\}}
  {\rLREb{\FormulaRefAuto{B48ContextAssumption}[thm],\rAI{13,11}}}
 \proofstepwidestar[2,3,5,9]{M,s[x\mapsto a]\models q}
  {\FormulaRefAuto{B48ConsequenceInstance}[thm]{2,14,12}}
 \proofstepwidestar[4,9]{(M,s\models q)\leftrightarrow(M,s[x\mapsto a]\models q)}
  {\FormulaRefAuto{B48FreshUpdate}[thm]{4,10}}
 \proofstepwidestar[2,3,4,5,9]{M,s\models q}
  {\rLREb{16,15}}
 \proofstepwidestar[1,2,3,4,5]{M,s\models q}
  {\rEE{8,9:17}}
 \proofstepwidestar[1,2,3,4]{(M,s\models \Gamma\cup\Delta)\to(M,s\models q)}
  {\rRI{5,18}}
 \proofstepwidestar[1,2,3,4]{\Gamma\cup\Delta\models_M q}
  {\FormulaRefAuto{B48SemanticConsequence}[def]{\rUI{19}}}
 \closeproofpartwide
\end{tabproofsplitwide}

\FormulaThmDeltaK[Semantische Gleichheitseinfuehrung]
{\Gamma\models_M v_x=v_x}
{B48RuleEqualityIntro}{
 \DeltaRow{Struktur}{\mathsf{Str}(M)}
 \DeltaRow{Annahmen}{\Gamma,\Delta,\Theta\subseteq\mathcal F}
 \DeltaRow{Formeln und Variablen}{p,q,r\in\mathcal F,\quad x,y,z\in\mathbb N}
}
\begin{tabproofsplitwide}
 \proofpartwide{Beweis}
 \proofstepwidestar[1]{M,s\models \Gamma}
  {\rA}
 \proofstepwidestar[]{s(x)=s(x)}
  {\rII}
 \proofstepwidestar[]{M,s\models v_x=v_x}
  {\FormulaRefAuto{B48Satisfaction}[def]{2}}
 \proofstepwidestar[]{(M,s\models \Gamma)\to(M,s\models v_x=v_x)}
  {\rRI{1,3}}
 \proofstepwidestar[]{\Gamma\models_M v_x=v_x}
  {\FormulaRefAuto{B48SemanticConsequence}[def]{\rUI{4}}}
 \closeproofpartwide
\end{tabproofsplitwide}

\FormulaThmDeltaK[Semantische Gleichheitsbeseitigung]
{\begin{gathered}
 \Gamma\models_M v_x=v_y,\quad\Delta\models_M p[x/z],\\
 \mathsf{Fr}(p,z,x),\quad\mathsf{Fr}(p,z,y)\\
 \vdash\ \Gamma\cup\Delta\models_M p[y/z]
 \end{gathered}}
{B48RuleEqualityElim}{
 \DeltaRow{Struktur}{\mathsf{Str}(M)}
 \DeltaRow{Annahmen}{\Gamma,\Delta,\Theta\subseteq\mathcal F}
 \DeltaRow{Formeln und Variablen}{p,q,r\in\mathcal F,\quad x,y,z\in\mathbb N}
}
\begin{tabproofsplitwide}
 \proofpartwide{Beweis}
 \proofstepwidestar[1]{\Gamma\models_M v_x=v_y}
  {\rA}
 \proofstepwidestar[2]{\Delta\models_M p[x/z]}
  {\rA}
 \proofstepwidestar[3]{\mathsf{Fr}(p,z,x)}
  {\rA}
 \proofstepwidestar[4]{\mathsf{Fr}(p,z,y)}
  {\rA}
 \proofstepwidestar[5]{M,s\models \Gamma\cup\Delta}
  {\rA}
 \proofstepwidestar[5]{(M,s\models \Gamma)\land(M,s\models \Delta)}
  {\rLREa{\FormulaRefAuto{B48ContextUnion}[thm],5}}
 \proofstepwidestar[1,5]{M,s\models v_x=v_y}
  {\FormulaRefAuto{B48ConsequenceInstance}[thm]{1,\rAEa{6}}}
 \proofstepwidestar[1,5]{s(x)=s(y)}
  {\FormulaRefAuto{B48Satisfaction}[def]{7}}
 \proofstepwidestar[2,5]{M,s\models p[x/z]}
  {\FormulaRefAuto{B48ConsequenceInstance}[thm]{2,\rAEb{6}}}
 \proofstepwidestar[3]{(M,s\models p[x/z])\leftrightarrow(M,s[z\mapsto s(x)]\models p)}
  {\FormulaRefAuto{B48SubstitutionLemma}[thm]{3}}
 \proofstepwidestar[2,3,5]{M,s[z\mapsto s(x)]\models p}
  {\rLREa{10,9}}
 \proofstepwidestar[1,2,3,5]{M,s[z\mapsto s(y)]\models p}
  {\rIE{8,11}}
 \proofstepwidestar[4]{(M,s\models p[y/z])\leftrightarrow(M,s[z\mapsto s(y)]\models p)}
  {\FormulaRefAuto{B48SubstitutionLemma}[thm]{4}}
 \proofstepwidestar[1,2,3,4,5]{M,s\models p[y/z]}
  {\rLREb{13,12}}
 \proofstepwidestar[1,2,3,4]{(M,s\models \Gamma\cup\Delta)\to(M,s\models p[y/z])}
  {\rRI{5,14}}
 \proofstepwidestar[1,2,3,4]{\Gamma\cup\Delta\models_M p[y/z]}
  {\FormulaRefAuto{B48SemanticConsequence}[def]{\rUI{15}}}
 \closeproofpartwide
\end{tabproofsplitwide}

\FormulaThmDeltaK[Semantischer Annahmenfall]
{p\in\Gamma\ \vdash\ \Gamma\models_M p}
{B48RuleAssumption}{
 \DeltaRow{Struktur}{\mathsf{Str}(M)}
 \DeltaRow{Annahmen}{\Gamma,\Delta,\Theta\subseteq\mathcal F}
 \DeltaRow{Formeln und Variablen}{p,q,r\in\mathcal F,\quad x,y,z\in\mathbb N}
}
\begin{tabproofsplitwide}
 \proofpartwide{Beweis}
 \proofstepwidestar[1]{p\in\Gamma}
  {\rA}
 \proofstepwidestar[2]{M,s\models \Gamma}
  {\rA}
 \proofstepwidestar[1,2]{M,s\models p}
  {\FormulaRefAuto{B48SemanticConsequence}[def]{2,1}}
 \proofstepwidestar[1]{(M,s\models \Gamma)\to(M,s\models p)}
  {\rRI{2,3}}
 \proofstepwidestar[1]{\Gamma\models_M p}
  {\FormulaRefAuto{B48SemanticConsequence}[def]{\rUI{4}}}
 \closeproofpartwide
\end{tabproofsplitwide}

\subsection{Induktion ueber die endliche Herleitung}

\FormulaDefDeltaK[Daten einer elementaren Beweistabelle]
{\begin{aligned}
 d&=((I_j,\theta_j,r_j,J_j))_{1\leq j\leq N},\\
 \Gamma_j&\coloneq\{\theta_k\mid k\in I_j\},\\
 H_d(n)&\coloneq
  \forall j\in\mathbb N\;
  ((1\leq j\leq N\land j\leq n)\to\Gamma_j\models_M\theta_j).
\end{aligned}}
{B48DerivationData}{
 \DeltaRow{Laenge}{N\in\mathbb N,\quad N\geq1}
 \DeltaRow{Formeln}{\theta_j\in\mathcal F}
 \DeltaRow{Annahmenindizes}{I_j\subseteq\{1,\ldots,j\}}
 \DeltaRow{Regel und fruehere Verweise}{r_j,\quad J_j}
}
Der Datensatz \(d\) ist eine \emph{gueltige} elementare Beweistabelle
nach Band~1: Eine Annahmezeile \(j\) hat \(I_j=\{j\}\).
Jeder sonstige Verweis in \(J_j\) zeigt auf eine Zeile mit Nummer
kleiner als \(j\). Die Formel und die Abhaengigkeitsmenge jeder
Schlusszeile entstehen nach ihrer angegebenen Regel.
Bei einer Entladung wird der Index der betreffenden Annahme aus
der Abhaengigkeitsmenge des Subbeweises entfernt; danach werden die
verbleibenden Mengen vereinigt.
Bei \(\forall I\) kommt die Variable in keiner offenen Annahme frei
vor. Bei \(\exists E\) kommt sie weder in der Schlussformel noch in
den nach der Entladung verbleibenden Annahmen frei vor.
Alle Einsetzungen erfuellen \(\mathsf{Fr}\).

Die \(\Gamma_j\) sind Mengen von Formeln, aber die Entladung wird
\emph{vorher an den Indizes} ausgefuehrt. Zwei verschiedene
Annahmezeilen mit derselben Formel duerfen deshalb nicht
miteinander verwechselt werden. Fuer eine entladene Annahme mit
Formel \(p\) gilt stets
\(\Gamma_{\rm Subbeweis}=\Gamma_{\rm Rest}\cup\{p\}\), auch wenn
\(p\) wegen einer anderen Annahme weiterhin in \(\Gamma_{\rm Rest}\)
liegt.

\FormulaThmDeltaK[Korrektheit jeder Zeile einer Beweistabelle]
{\forall j\in\{1,\ldots,N\}\;(\Gamma_j\models_M\theta_j)}
{B48DerivationSoundness}{
 \DeltaRow{Struktur}{\mathsf{Str}(M)}
 \DeltaRow{Gueltige elementare Beweistabelle}{d\text{ nach }
                          \FormulaRefAuto{B48DerivationData}[def]}
}
Wir beweisen \(H_d(n)\) durch Induktion ueber \(n\).
Fuer \(n=0\) ist der Bereich der zu pruefenden Zeilen leer.
Sei \(H_d(n)\) gegeben. Fuer \(n+1>N\) kommt keine Zeile hinzu.
Andernfalls betrachten wir die Regel der Zeile \(n+1\).
Alle unten als semantische Folgerungen aufgefuehrten Premissen
folgen aus \(H_d(n)\), weil die referenzierten Zeilen frueher
liegen. Die Einsetzungs- und Eigenvariablenbedingungen folgen
aus der Gueltigkeit von \(d\). Die Mengen
\(\Gamma,\Delta,\Theta\) sind jeweils die verbleibenden
Annahmenmengen der betreffenden Praemissen beziehungsweise
Subbeweise.

\begin{longtable}{@{}p{0.13\linewidth}p{0.56\linewidth}p{0.24\linewidth}@{}}
\toprule
Regel & Angewendeter semantischer Schluss & Nachweis\\
\midrule\endhead
\(\RuleRef{A}{A}\) &
 \(p\in\Gamma\ \vdash\ \Gamma\models_M p\) &
 \FormulaRefAuto{B48RuleAssumption}[thm]\\[5pt]
\(\RuleRef{AI}{\land I}\) &
 \(\begin{gathered}\Gamma\models_M p,\ \Delta\models_M q\\
 \vdash\Gamma\cup\Delta\models_M p\land q\end{gathered}\) &
 \FormulaRefAuto{B48RuleAndIntro}[thm]\\[5pt]
\(\RuleRef{AE1}{\land E_1}\) &
 \(\Gamma\models_M p\land q\ \vdash\Gamma\models_M p\) &
 \FormulaRefAuto{B48RuleAndLeft}[thm]\\[5pt]
\(\RuleRef{AE2}{\land E_2}\) &
 \(\Gamma\models_M p\land q\ \vdash\Gamma\models_M q\) &
 \FormulaRefAuto{B48RuleAndRight}[thm]\\[5pt]
\(\RuleRef{OI1}{\lor I_1}\) &
 \(\Gamma\models_M p\ \vdash\Gamma\models_M p\lor q\) &
 \FormulaRefAuto{B48RuleOrLeft}[thm]\\[5pt]
\(\RuleRef{OI2}{\lor I_2}\) &
 \(\Gamma\models_M q\ \vdash\Gamma\models_M p\lor q\) &
 \FormulaRefAuto{B48RuleOrRight}[thm]\\[5pt]
\(\RuleRef{OE}{\lor E}\) &
 \(\begin{gathered}\Gamma\models_M p\lor q,\\
 \Delta\cup\{p\}\models_M r,\ \Theta\cup\{q\}\models_M r\\
 \vdash(\Gamma\cup\Delta)\cup\Theta\models_M r\end{gathered}\) &
 \FormulaRefAuto{B48RuleOrElim}[thm]\\[5pt]
\(\RuleRef{RI}{\to I}\) &
 \(\Gamma\cup\{p\}\models_M q\ \vdash\Gamma\models_M p\to q\) &
 \FormulaRefAuto{B48RuleImpIntro}[thm]\\[5pt]
\(\RuleRef{RE}{\to E}\) &
 \(\begin{gathered}\Gamma\models_M p\to q,\ \Delta\models_M p\\
 \vdash\Gamma\cup\Delta\models_M q\end{gathered}\) &
 \FormulaRefAuto{B48RuleImpElim}[thm]\\[5pt]
\(\RuleRef{LRI}{\leftrightarrow I}\) &
 \(\begin{gathered}\Gamma\models_M p\to q,\ \Delta\models_M q\to p\\
 \vdash\Gamma\cup\Delta\models_M p\leftrightarrow q\end{gathered}\) &
 \FormulaRefAuto{B48RuleIffIntro}[thm]\\[5pt]
\(\RuleRef{LRE1}{\leftrightarrow E_1}\) &
 \(\Gamma\models_M p\leftrightarrow q\ \vdash\Gamma\models_M p\to q\) &
 \FormulaRefAuto{B48RuleIffLeft}[thm]\\[5pt]
\(\RuleRef{LRE2}{\leftrightarrow E_2}\) &
 \(\Gamma\models_M p\leftrightarrow q\ \vdash\Gamma\models_M q\to p\) &
 \FormulaRefAuto{B48RuleIffRight}[thm]\\[5pt]
\(\RuleRef{BI}{\bot I}\) &
 \(\begin{gathered}\Gamma\models_M p,\ \Delta\models_M\neg p\\
 \vdash\Gamma\cup\Delta\models_M\bot\end{gathered}\) &
 \FormulaRefAuto{B48RuleBottom}[thm]\\[5pt]
\(\RuleRef{CI}{\neg I}\) &
 \(\Gamma\cup\{p\}\models_M\bot\ \vdash\Gamma\models_M\neg p\) &
 \FormulaRefAuto{B48RuleNegIntro}[thm]\\[5pt]
\(\RuleRef{CE}{\neg E}\) &
 \(\Gamma\cup\{\neg p\}\models_M\bot\ \vdash\Gamma\models_M p\) &
 \FormulaRefAuto{B48RuleClassical}[thm]\\[5pt]
\(\RuleRef{UI}{\forall I}\) &
 \(\begin{gathered}\Gamma\models_M p,\ x\mathrel{\#}\Gamma\\
 \vdash\Gamma\models_M\forall x\,p\end{gathered}\) &
 \FormulaRefAuto{B48RuleAllIntro}[thm]\\[5pt]
\(\RuleRef{UE}{\forall E}\) &
 \(\begin{gathered}\Gamma\models_M\forall x\,p,\ \mathsf{Fr}(p,x,y)\\
 \vdash\Gamma\models_M p[y/x]\end{gathered}\) &
 \FormulaRefAuto{B48RuleAllElim}[thm]\\[5pt]
\(\RuleRef{EI}{\exists I}\) &
 \(\begin{gathered}\Gamma\models_M p[y/x],\ \mathsf{Fr}(p,x,y)\\
 \vdash\Gamma\models_M\exists x\,p\end{gathered}\) &
 \FormulaRefAuto{B48RuleExistsIntro}[thm]\\[5pt]
\(\RuleRef{EE}{\exists E}\) &
 \(\begin{gathered}\Gamma\models_M\exists x\,p,\ \Delta\cup\{p\}\models_M q,\\
 x\mathrel{\#}\Delta,\ x\notin\operatorname{FV}(q)\\
 \vdash\Gamma\cup\Delta\models_M q\end{gathered}\) &
 \FormulaRefAuto{B48RuleExistsElim}[thm]\\[5pt]
\(\RuleRef{II}{=I}\) &
 \(\Gamma\models_M v_x=v_x\) &
 \FormulaRefAuto{B48RuleEqualityIntro}[thm]\\[5pt]
\(\RuleRef{IE}{=E}\) &
 \(\begin{gathered}\Gamma\models_M v_x=v_y,\ \Delta\models_M p[x/z],\\
 \mathsf{Fr}(p,z,x),\ \mathsf{Fr}(p,z,y)\\
 \vdash\Gamma\cup\Delta\models_M p[y/z]\end{gathered}\) &
 \FormulaRefAuto{B48RuleEqualityElim}[thm]\\
\bottomrule
\end{longtable}

Damit ist in jedem moeglichen Regelfall
\(\Gamma_{n+1}\models_M\theta_{n+1}\) nachgewiesen.
Fuer \(j\leq n\) gilt die Behauptung weiter nach \(H_d(n)\);
fuer \(j=n+1\) gilt der gerade bewiesene Fall. Also gilt
\(H_d(n+1)\). Die Induktion
\FormulaRefAuto{PeanoInductionRuleAddOne}[thm] liefert
\(\forall n\in\mathbb N\;H_d(n)\). Einsetzen von \(n=N\) ist
genau die behauptete Korrektheit aller Zeilen.

\noindent
Die aufgefuehrten 21 Faelle decken alle elementaren Regeln aus
Band~1 ab. Die dortigen iterierten Gleichheits-, Konjunktions-
und Existenznotationen bezeichnen endliche Folgen solcher
Regelschritte. Der Satz setzt keine Korrektheit ungepruefter
Bibliotheksverweise voraus: \(\operatorname{Prf}_1\) bezeichnet hier
die elementare Herleitung, in der verwendete Abkuerzungen bereits
aufgeloest sind. Die fuer Henkin benoetigten allgemeinen
Umformungssaetze ueber Beweisobjekte bleiben eine eigene Aufgabe.

\FormulaThmDeltaK[Ein Modell erfuellt jede Teilmenge seiner Theorie]
{M\models T,\quad\Gamma\subseteq T\ \vdash\ M,s\models\Gamma}
{B48ModelContext}{
 \DeltaRow{Struktur und Belegung}{\mathsf{Str}(M),\quad s\in A_D}
 \DeltaRow{Theorie und Annahmen}{T\subseteq\operatorname{Sent}(\mathcal L_\in),
                               \quad\Gamma\subseteq\mathcal F}
}
\begin{tabproof}
 \proofstep{1}{M\models T}{\rA}
 \proofstep{2}{\Gamma\subseteq T}{\rA}
 \proofstep{3}{p\in\Gamma}{\rA}
 \proofstep{2,3}{p\in T}
  {\FormulaRefAuto{A\subseteq B,\ x\in A\vdash x\in B}[thm]{2,3}}
 \proofstep{1,2,3}{M\models p}
  {\FormulaRefAuto{B48SentenceSatisfaction}[def]{1,4}}
 \proofstep{}{s\colon\operatorname{Var}\to D}
  {\FormulaRefAuto{B48AssignmentCriterion}[thm]}
 \proofstep{1,2,3}{M,s\models p}
  {\FormulaRefAuto{B48SentenceInstance}[thm]{5,6}}
 \proofstep{1,2}{M,s\models\Gamma}
  {\FormulaRefAuto{B48SemanticConsequence}[def]{\rUI{\rRI{3,7}}}}
\end{tabproof}

\FormulaThmDeltaK[Korrektheit der Herleitung]
{T\vdash\varphi,\quad M\models T\ \vdash\ M\models\varphi}
{B48Soundness}{
 \DeltaRow{Struktur}{\mathsf{Str}(M)}
 \DeltaRow{Theorie}{T\subseteq\operatorname{Sent}(\mathcal L_\in)}
 \DeltaRow{Satz}{\varphi\in\operatorname{Sent}(\mathcal L_\in)}
}
Im Zeugenbeweis ist \(d\) frisch. Seine Laenge ist \(N\), seine
letzte Formel ist \(\varphi\), und seine offenen Annahmen
\(\Gamma_N\) liegen in \(T\). Die Belegung \(s\) wird beliebig
gewaehlt und vor der Entladung von \(d\) wieder gebunden.
\begin{tabproofsplitwide}
 \proofpartwide{Schluss aus der letzten Beweiszeile}
 \proofstepwidestar[1]{T\vdash\varphi}{\rA}
 \proofstepwidestar[2]{M\models T}{\rA}
 \proofstepwidestar[1]{\exists d\;\operatorname{Prf}_1(d,T,\varphi)}
  {\FormulaRefAuto{B48Syntax}[def]{1}}
 \proofstepwidestar[4]{\operatorname{Prf}_1(d,T,\varphi)}{\rA}
 \proofstepwidestar[4]{\Gamma_N\models_M\varphi}
  {\FormulaRefAuto{B48DerivationSoundness}[thm]{4},
   \FormulaRefAuto{B48DerivationData}[def]}
 \proofstepwidestar[4]{\Gamma_N\subseteq T}
  {\FormulaRefAuto{B48Syntax}[def]{4}}
 \proofstepwidestar[7]{s\in A_D}{\rA}
 \proofstepwidestar[2,4,7]{M,s\models\Gamma_N}
  {\FormulaRefAuto{B48ModelContext}[thm]{2,6,7}}
 \proofstepwidestar[2,4,7]{M,s\models\varphi}
  {\FormulaRefAuto{B48ConsequenceInstance}[thm]{5,8}}
 \proofstepwidestar[2,4]{\forall s\in A_D\;(M,s\models\varphi)}
  {\rUI{\rRI{7,9}}}
 \proofstepwidestar[]{\exists a\;(a\in D)}
  {\FormulaRefAuto{B48StrDomainAxiom}[ax]}
 \proofstepwidestar[2,4]{M\models\varphi}
  {\FormulaRefAuto{B48SentenceSatisfaction}[def]{11,10},
   \FormulaRefAuto{B48AssignmentCriterion}[thm]}
 \proofstepwidestar[1,2]{M\models\varphi}{\rEE{3,4:12}}
 \closeproofpartwide
\end{tabproofsplitwide}

